% ===================================================================
%  BẢNG Số 16 — Cửa hàng lớn. --- Chợ phiên. Đồ chơi.
%
%  Traduction, faite en 2026, en vietnamien standard d'aujourd'hui,
%  sur l'ido de texto/io. Regles de la colonne : voir 10-bang-01.tex.
%  Dernier tableau du livret.
%
%  LES ANIMAUX EN CARTON DU RAYON DES JOUETS SONT LE MEILLEUR RESUME
%  DE TOUTE LA COLONNE. Ceux que ce pays connait portent leur nom
%  vietnamien : khỉ le singe, cáo le renard, cá sấu le crocodile, hổ
%  le tigre, voi l'elephant, sói le loup, chuột le rat, rắn le
%  serpent, rùa la tortue. Ceux qu'il ne connait pas portent un nom
%  sino-vietnamien, forge par description : « tê giác » la corne de
%  rhinoceros, « hà mã » le cheval de riviere pour l'hippopotame,
%  « lạc đà » le chameau, « hươu cao cổ » le cerf a long cou pour la
%  girafe, « đà điểu » l'autruche. Et un seul est reste francais,
%  « leo-pa », parce qu'il est entre par le cirque.
%
%  LE MEME PARTAGE VAUT POUR LA BATAILLE DE L'ALINEA 5, qui est le
%  passage le plus sino-vietnamien du livret entier : quân đội
%  l'armee, pháo đài la forteresse, bao vây le siege, tấn công
%  l'assaut, đầu hàng la reddition, hòa ước le traite de paix. Le
%  vocabulaire militaire est venu des annales chinoises, et il n'a
%  pas change depuis.
%
%  LE TABLEAU D'ASTRONOMIE FERME LE LIVRET SUR LE MEME CONSTAT QUE LE
%  TABLEAU 1 : sao chổi la comete, nhật thực l'eclipse de soleil,
%  nguyệt thực celle de lune, hành tinh la planete, bán cầu
%  l'hemisphere, hệ mặt trời le systeme solaire — tout est chinois,
%  parce que tout s'apprend a l'ecole. Cent six pages plus tot, la
%  meme classe portait les memes mots.
%
%  LES ETRENNES DU QUEBEC ETAIENT DEVENUES « CADEAUX DU JOUR DE
%  L'AN » ; ici elles sont « quà Tết », les cadeaux du Tet, et le mot
%  est exact : le nouvel an est ici la fete des cadeaux, comme il
%  l'etait dans la France de 1926. L'arbre de Noel du renvoi 81 reste
%  ce qu'il est, « cây Nô-en », arrive avec la mission.
% ===================================================================

%%K t16-apar-1 apar
\VUsaut{4.0mm}
\VUtitre{15.0pt}{60}{BẢNG Số 16}
\VUsaut{2.0mm}
\VUpk{13.2pt}{40}{Cửa hàng lớn. --- Chợ phiên.}
\VUpk{13.2pt}{40}{Đồ chơi.}
\VUsaut{3.4mm}

%%K t16-01-1 p
1. --- Năm mới sắp đến. Các cửa hàng lớn đã sửa soạn những quầy bày \VUgras{đồ chơi}
\textsuperscript{(1)} và \VUgras{quà Tết}.

%%K t16-01-2 p
Em xin ông nội dẫn em đến chợ phiên để xem cuộc trưng bày đẹp đẽ ấy, và ông đã nhận lời.

%%K t16-02-1 p
2. --- Trước hết chúng em dừng lại trước những tấm bảng treo vũ khí đẹp, trên đó có một
khẩu \VUgras{súng}, một chiếc \VUgras{mũ lưỡi trai cứng}, một thanh \VUgras{gươm}, một
chiếc \VUgras{thắt lưng đeo kiếm} và một đôi \VUgras{cầu vai}, hoặc một chiếc \VUgras{mũ
sắt}, một tấm \VUgras{áo giáp} \textsuperscript{(3)} và một khẩu \VUgras{súng lục}
\textsuperscript{(4)}. Tất cả những thứ ấy làm em thích thú hơn nhiều so với những chiếc
\VUgras{đèn pha} \textsuperscript{(2)}.

%%K t16-tit-1 sub
\VUsaut{3.0mm}
\VUpk{10.2pt}{40}{Những thứ đẹp mắt.}
\VUsaut{2.6mm}

%%K t16-03-1 p
3. --- Cùng quầy ấy có một \VUgras{người lính gác} \textsuperscript{(5)} đứng trong
\VUgras{chòi canh} của mình; anh ta như đang canh trước cả một vườn thú bằng bìa. Ở đó có
biết bao là con vật : một con \VUgras{vẹt} \textsuperscript{(6)} trên cây đậu, một con
\VUgras{tê giác} \textsuperscript{(7)} có sừng trên mũi, một con \VUgras{tuần lộc}
\textsuperscript{(8)}, một con \VUgras{đà điểu} \textsuperscript{(9)}, một con
\VUgras{khỉ} \textsuperscript{(10)} đang cắn quả hồ đào, một con \VUgras{cáo}
\textsuperscript{(11)} tha con gà mái, một con \VUgras{cá sấu} \textsuperscript{(12)} há
đôi hàm khổng lồ, một con \VUgras{lạc đà} \textsuperscript{(13)}, một con \VUgras{chó săn
thỏ} \textsuperscript{(14)} thanh nhã, một con \VUgras{báo đen} \textsuperscript{(15)},
rồi hai con vật em không biết và người ta bảo là con \VUgras{chồn}
\textsuperscript{(16)} và con \VUgras{linh cẩu} \textsuperscript{(17)}. Ở đó cũng có một
con \VUgras{leo-pa} \textsuperscript{(18)} và một con \VUgras{sói}
\textsuperscript{(21)}; ngay gần đó là những con \VUgras{chuột} \textsuperscript{(19)}
nhỏ xinh và những con \VUgras{chuột nhắt} \textsuperscript{(20)} tí hon bằng cao su. Sau
cùng, một con \VUgras{hà mã} \textsuperscript{(22)} và một con \VUgras{lạc đà một bướu}
\textsuperscript{(23)} khép lại bộ sưu tập. Em đã ở đó hàng giờ nếu bên cạnh không có
một trại lính lộng lẫy đầy \VUgras{lính chì} \textsuperscript{(29)} thu hút sự chú ý của
em.

%%K t16-tit-2 sub
\VUsaut{4.0mm}
\VUpk{10.2pt}{40}{Lính và chiến tranh.}
\VUsaut{2.8mm}

%%K t16-04-1 p
4. --- Ông nội em, một \VUgras{thương binh} \textsuperscript{(24)} đã phải rời quân ngũ
vì một \VUgras{vết thương} mang lại cho ông tấm \VUgras{huy chương quân công}, rất thích
kể cho em nghe những \VUgras{chiến dịch} mà ông đã dự. Ông sung sướng giảng cho em các
thế trận khác nhau của đạo quân tí hon ấy. Một nhóm lính thật đang thăm chợ phiên — một
\VUgras{lính công binh} \textsuperscript{(25)}, một \VUgras{lính sơn cước}
\textsuperscript{(26)} đội \VUgras{mũ bê-rê}, một \VUgras{trung sĩ nhất}
\textsuperscript{(27)} bộ binh và một \VUgras{trung sĩ pháo binh}
\textsuperscript{(28)} với chiếc \VUgras{lon} vàng trên tay áo — dừng lại gần chúng em để
nghe ông giảng.

%%K t16-05-1 p
5. --- Ông nội em, rất tự hào vì có một cử tọa lừng lẫy như thế, liền ứng khẩu vạch ra
cả một kế hoạch trận đánh.

%%K t16-05-2 p
Thành phố có công sự, với \VUgras{tường thành} và \VUgras{hào}, bị \VUgras{quân địch}
\VUgras{bao vây}, và sau một cuộc \VUgras{bắn phá} lâu dài, quân địch đã sẵn sàng
\VUgras{tấn công}. Nhưng \VUgras{quân bị vây}, bị cái đói thúc ép, đã thử \VUgras{phá
vây} đánh ra \VUgras{trại} của \VUgras{quân vây}. Sau khi đẩy lui \VUgras{cuộc tấn công}
bằng một \VUgras{trận đánh} dai dẳng quanh \VUgras{doanh trại}, \VUgras{quân thắng trận}
tung \VUgras{kỵ binh} truy đuổi những kẻ tấn công \VUgras{đã bại}. Những người này
\VUgras{rút lui}, để lại trên \VUgras{chiến trường} đầy \VUgras{người chết} và
\VUgras{người bị thương}. \VUgras{Người khiêng cáng} mang những người sau cùng ấy về
\VUgras{trạm cứu thương} để \VUgras{bác sĩ phẫu thuật} chăm sóc.

%%K t16-05-3 p
Dù vài \VUgras{tiểu đoàn} kháng cự anh dũng bằng cách dàn thành \VUgras{ô vuông}, cuộc
\VUgras{bại trận} vẫn hoàn toàn, phần còn lại của quân đội \VUgras{bỏ chạy}, để lại rất
nhiều \VUgras{tù binh}. Quân địch sẽ hoàn tất \VUgras{chiến thắng} của mình bằng cách
chiếm thành phố. Có lẽ đó sẽ là kết thúc của chiến dịch và, sau một cuộc \VUgras{đình
chiến}, sẽ ký được một \VUgras{hòa ước}.

%%K t16-tit-3 sub
\VUsaut{4.4mm}
\VUpk{10.2pt}{40}{Quà Tết.}
\VUsaut{2.8mm}

%%K t16-06-1 p
6. --- Mấy người lính trẻ, cười vì câu chuyện sinh động của người cựu binh, hỏi ông thêm
vài điều nữa. Còn em, em đi vòng quanh gian hàng để xem một cái \VUgras{nỏ}
\textsuperscript{(32)} đẹp và một cây \VUgras{cung} \textsuperscript{(33)} với
\VUgras{mũi tên} \textsuperscript{(31)} của nó. Ở đó em tìm thấy một bộ sưu tập thú vật
khác : hai con \VUgras{chim hót} \textsuperscript{(34)}, một con \VUgras{họa mi} chạy máy
và một con \VUgras{chim chích} hót vang cả cổ, và một loạt con vật kỳ lạ : một con
\VUgras{dơi} \textsuperscript{(35)}, một con \VUgras{trăn} \textsuperscript{(36)}, một
con \VUgras{rùa} \textsuperscript{(37)}, một con \VUgras{hải cẩu} \textsuperscript{(38)},
một con \VUgras{cá voi} \textsuperscript{(39)} và một con \VUgras{cá mập}
\textsuperscript{(40)} bằng bong bóng.

%%K t16-07-1 p
7. --- Em không dừng lại lâu để ngắm chúng, vì em đã thấy thứ em hằng mơ ước : một
\VUgras{rạp múa rối} \textsuperscript{(41)} tuyệt đẹp với những tấm \VUgras{phông} lộng
lẫy và một tấm \VUgras{màn} đỏ mê hồn. Trên \VUgras{sân khấu}, hai diễn viên như đang
đóng một \VUgras{vai} trong một \VUgras{vở kịch} rất hấp dẫn, vì những \VUgras{khán giả}
tí hon bằng bìa ngồi trong các lô hay trên những chiếc \VUgras{ghế bành} và ở
\VUgras{tầng trệt} phía sau \VUgras{người nhạc trưởng} đều vỗ tay rất nồng nhiệt.

%%K t16-07-2 p
Tiếc thay, rạp hát ấy đắt quá.

%%K t16-08-1 p
8. --- Vả lại, chọn một món đồ chơi thật khó, vì các tủ kính chất đầy đủ loại : một
\VUgras{chú hề Arlequin} \textsuperscript{(42)}, một \VUgras{chú Pôlisinen}
\textsuperscript{(43)}, một \VUgras{chú Pi-e-rô} \textsuperscript{(44)}, những con
\VUgras{cú} \textsuperscript{(45)} vỗ cánh, những con \VUgras{sáo}
\textsuperscript{(46)} biết huýt, những con \VUgras{chó xù} biết sủa, một chiếc
\VUgras{máy hát} \textsuperscript{(47)} và những \VUgras{trò chơi trí tuệ}. Một trong
những trò ấy làm em rất thích : nó bày ra một \VUgras{trận thủy chiến}
\textsuperscript{(48)}, trong đó một chiếc \VUgras{tàu} bị \VUgras{bọn cướp biển} tấn
công gần một vùng đất trồng \VUgras{dừa}.

%%K t16-noto-1 noto
\VUnotes{143.2mm}{%
\textit{(*) Xem bảng số 6.}%
}

%%K t16-09-1 p
9. --- Em có thể thong thả ngắm tất cả những điều kỳ diệu ấy, vì trong cửa hàng chỉ có
ít người : vài người đi dạo, một \VUgras{lính kỵ binh} \textsuperscript{(49)} và một
\VUgras{người thu tiền} \textsuperscript{(50)} đang trình một \VUgras{hối phiếu} ở
\VUgras{quầy thu ngân chính} \textsuperscript{(51)}.

%%K t16-10-1 p
10. --- Em hỏi ông nội xem những cuốn sổ đặt trên các tấm ván sau lưng \VUgras{chị thu
ngân} dùng để làm gì; ông trả lời rằng đó là những \VUgras{sổ kế toán}.

%%K t16-tit-4 sub
\VUsaut{3.6mm}
\VUpk{10.2pt}{40}{Ngắm các tủ kính.}
\VUsaut{2.8mm}

%%K t16-11-1 p
11. --- Rời quầy đồ chơi, chúng em sang quầy đồ yên cương để nhìn qua những chiếc
\VUgras{yên} \textsuperscript{(53)}, \VUgras{bàn đạp} \textsuperscript{(54)},
\VUgras{dây cương} \textsuperscript{(55)}, \VUgras{hàm thiếc} \textsuperscript{(56)} và
\VUgras{roi ngựa}. Trong lúc \VUgras{trưởng quầy} \textsuperscript{(57)} giải thích vài
điều cho ông nội em, em đi xem quầy bày \VUgras{đồ sứ} và \VUgras{đồ pha lê}
\textsuperscript{(58)}, nơi có những chiếc \VUgras{bát đựng xà lách} đẹp và những
\VUgras{ấm trà} \textsuperscript{(59)} tinh xảo.

%%K t16-12-1 p
12. --- Để lên tầng một, chúng em đi qua sau tủ kính bày \VUgras{áo nịt}
\textsuperscript{(60)}, cạnh quầy hàng dệt kim, nơi bày những chiếc \VUgras{quần lót}
\textsuperscript{(63)} rất đẹp. Gần đó, một \VUgras{cô bán hàng} \textsuperscript{(61)}
đang mời một \VUgras{bà khách} \textsuperscript{(62)} chọn những tấm \VUgras{lụa}
\textsuperscript{(64)} đẹp.

%%K t16-12-2 p
Khi lên cầu thang, em nhận thấy cửa hàng được trang trí bằng những chiếc \VUgras{đèn
lồng} \textsuperscript{(65)} Nhật Bản, những chiếc \VUgras{ô} \textsuperscript{(66)} và
những cây \VUgras{cọ} \textsuperscript{(70)} khổng lồ.

%%K t16-13-1 p
13. --- Ở tầng một không có gì nhiều để xem : chỉ có đồ đạc (*), những chiếc \VUgras{két
sắt} \textsuperscript{(67)}, những chiếc \VUgras{tủ com-mốt} \textsuperscript{(68)} và
những chiếc \VUgras{xe nôi} \textsuperscript{(69)}. Nhưng toàn cảnh cửa hàng thì rất
\VUgras{vui mắt}.

%%K t16-14-1 p
14. --- Dưới chân chúng em, em thấy các nhạc cụ : những cây \VUgras{đàn hạc}
\textsuperscript{(71)}, \VUgras{đàn ghi-ta} \textsuperscript{(72)}, \VUgras{đàn măng-đô-lin}
\textsuperscript{(73)}, những chiếc \VUgras{kèn cor-nê} \textsuperscript{(74)},
\VUgras{kèn cla-ri-nét} \textsuperscript{(75)}, những cây vĩ cầm với \VUgras{ác-sê}
\textsuperscript{(76)} của chúng và cả một chiếc \VUgras{trống lớn}
\textsuperscript{(77)}. Xa hơn một chút, ở quầy văn phòng phẩm, một \VUgras{sinh viên}
\textsuperscript{(78)} đang mặc cả một chiếc cặp với \VUgras{người bán hàng}
\textsuperscript{(79)}. Gần họ, em thấy một cuốn \VUgras{sách vỡ lòng}
\textsuperscript{(80)} đẹp.

%%K t16-14-2 p
Giữa cửa hàng người ta đã dựng một \VUgras{cây Nô-en} \textsuperscript{(81)} lớn treo đầy
nến, đồ trang trí nhỏ và đủ loại đồ chơi : những con \VUgras{ngựa gỗ}
\textsuperscript{(82)}, những con \VUgras{búp bê} \textsuperscript{(83)}, v.v.

%%K t16-14-3 p
\VUgras{Quầy nước hoa} \textsuperscript{(84)}, với những thứ \VUgras{thuốc đánh răng} và
\VUgras{nước hoa} khác nhau, tỏa một mùi thơm dễ chịu bay lên tận chỗ chúng em.

%%K t16-tit-5 sub
\VUsaut{3.4mm}
\VUpk{10.2pt}{40}{Chúng em mua một thứ.}
\VUsaut{2.8mm}

%%K t16-15-1 p
15. --- Vì ông nội em bắt đầu thấy lâu, chúng em xuống bằng cầu thang lớn. Ông dừng lại
một lát trước một \VUgras{bảng thiên văn} \textsuperscript{(85)} vẽ, ông bảo em, những
\VUgras{sao chổi} \textsuperscript{(a)}, những \VUgras{lần nhật thực và nguyệt thực}
\textsuperscript{(b)}, một hoa gió với \VUgras{bốn phương chính} \textsuperscript{(c)}
\VUgras{(bắc, nam, đông và tây)}, một bản đồ hai \VUgras{bán cầu}
\textsuperscript{(d)}, các \VUgras{hành tinh} \textsuperscript{(e)} và \VUgras{hệ mặt
trời} \textsuperscript{(f)}. Em chẳng hiểu gì cả, và em thích thú hơn nhiều với bộ chế
phục viền lon của một \VUgras{người giao hàng} \textsuperscript{(86)} đi qua, và với
những chiếc \VUgras{quạt} \textsuperscript{(87)} đẹp ở gần em, cạnh những chiếc
\VUgras{ví tiền} \textsuperscript{(88)} và những chiếc \VUgras{ví da}
\textsuperscript{(89)}.

%%K t16-16-1 p
16. --- Trên quầy hàng em cũng thán phục những chiếc \VUgras{lông cắm mũ}
\textsuperscript{(90)} và những chiếc \VUgras{khóa thắt lưng} \textsuperscript{(91)},
cùng những bông \VUgras{hoa giả} \textsuperscript{(94)} đẹp đẽ xếp duyên dáng trong các
giỏ. Những bông \VUgras{hoa tím}, \VUgras{hoa quế trúc}, \VUgras{hoa lan dạ hương}
\textsuperscript{(95)}, \VUgras{hoa phong lữ} \textsuperscript{(96)}, \VUgras{hoa huệ
tây} \textsuperscript{(97)} và \VUgras{hoa sen cạn} \textsuperscript{(98)} đều được làm
giống hệt.

%%K t16-tit-6 sub
\VUsaut{4.4mm}
\VUpk{10.2pt}{40}{Món quà của ông nội.}
\VUsaut{2.8mm}

%%K t16-17-1 p
17. --- Em xin ông nội mua cho em một trong những chiếc \VUgras{bàn chải đánh răng}
\textsuperscript{(92)} đựng trong một cái hộp, cạnh những chiếc \VUgras{trâm cài tóc}
\textsuperscript{(93)}. Ông nhận lời và em tự mình đi trả tiền món hàng ở quầy thu ngân,
gần đó người ta đang chuẩn bị một \VUgras{kiện hàng} \textsuperscript{(100)} lớn cho một
\VUgras{người khách} \textsuperscript{(99)}.

%%K t16-18-1 p
18. --- Bỗng có một chút lộn xộn giữa những người đang đứng gần quầy \VUgras{đồ đồng mỹ
nghệ} \textsuperscript{(101)}. Một \VUgras{tên trộm} \textsuperscript{(102)} vừa bị một
\VUgras{người kiểm soát} \textsuperscript{(103)} bắt quả tang lúc hắn đang định móc túi
một người. Người ta đã cho gọi cảnh sát đến \VUgras{bắt} hắn, và đúng lúc chúng em ra
về, kẻ phạm tội bị dẫn vào nhà giam.
\VUsaut{6.0mm}
\VUfilet{22mm}
